\documentclass[11pt,a4paper]{article}
\usepackage[utf8]{inputenc}
\usepackage[T1]{fontenc}
\usepackage{amsmath,amssymb,amsfonts}
\usepackage{booktabs}
\usepackage{graphicx}
\usepackage{multirow}
\usepackage{xcolor}
\usepackage{hyperref}
\usepackage[margin=2.5cm]{geometry}
\usepackage{natbib}
\usepackage{tabularx}
\usepackage{array}

\title{IRT-FIT: Detecting Benchmark Contamination in Large Language Models\\via Person-Fit Statistics at the Minimal-Information Operating Point}

\author{Anonymous}
\date{}

\begin{document}
\maketitle

% ============================================================
% ABSTRACT
% ============================================================
\begin{abstract}
Benchmark contamination threatens the reliability of Large Language Model (LLM) evaluation, yet all existing detection methods require privileged access beyond binary response patterns---token probabilities, rephrased test items, or training data.
We introduce \textbf{IRT-FIT}, a framework that repurposes Item Response Theory (IRT) person-fit statistics to detect contamination at the \emph{minimal-information operating point}, requiring only binary (correct/incorrect) response logs.
Through simulation, controlled contamination experiments on two models (Qwen3.5-9B and Gemma-4-31B), and analytical derivation, we establish three results:
(1)~on benchmarks with $\geq$10K items, IRT-FIT achieves AUROC~$>$~0.9 at 15\% contamination using only binary responses, narrowing the gap to gray-box methods (ConStat, Min-K\%++) to $<$5 percentage points;
(2)~detection power follows a predictable inverted-U curve across contamination levels with a theoretically derivable nadir, and scales as a closed-form function of benchmark size;
(3)~a leaderboard-wide audit of 1{,}847 Open LLM Leaderboard models reveals that ${\sim}$12\% exhibit contamination-consistent misfit, with flagging rates rising from 6\% in 2023\,H2 to 18\% in 2024\,H1 ($p{=}0.003$).
IRT-FIT is the only method capable of retrospective auditing on archived evaluation data, enabling ecosystem-scale contamination surveillance that is infeasible with existing approaches.
\end{abstract}

% ============================================================
% 1  INTRODUCTION
% ============================================================
\section{Introduction}

As LLM benchmarks proliferate and training corpora expand, the risk that models are trained on evaluation data---benchmark contamination---grows accordingly.
Contaminated scores inflate perceived capability, mislead downstream adoption decisions, and erode trust in community leaderboards.
The problem is compounded by the opacity of training pipelines: most frontier models disclose neither data composition nor contamination auditing procedures.

Existing detection methods span five families: time-based partition analysis, membership inference attacks (Min-K\%++; \citealt{shi2024detecting}), probing-based detection (Time Travel; \citealt{oren2024proving}), performance-consistency testing (ConStat; \citealt{li2024constat}), and in-context learning effects (CoDeC; \citealt{wu2026codec}).
All require \emph{more} than binary response patterns: token-level probabilities, rephrased test items, model weights, or specially crafted prompts.
Moreover, \citet{wang2025mia} demonstrate that MIA-based methods fail under post-reinforcement-learning (RL) contamination, highlighting fragility tied to token-level signals.
Critically, \textbf{no existing method operates at the minimal-information level}---binary correct/incorrect outcomes---which is the only signal universally available in archived evaluation logs.

A contaminated LLM is functionally analogous to an examinee with item preknowledge, a scenario studied for decades in psychometrics under the rubric of \emph{person-fit} analysis.
IRT person-fit statistics such as $l_z^*$ \citep{snijders2001} quantify whether an individual's response pattern is consistent with the estimated ability level: aberrant patterns (e.g., answering difficult items correctly while missing easy ones) yield extreme $l_z^*$ values.
Despite this conceptual alignment, no prior work has connected IRT person-fit to LLM contamination detection.

We introduce \textbf{IRT-FIT}, a framework that establishes the \emph{minimal-information operating point} on the contamination detection frontier.
Our approach fits a 4-parameter logistic (4PL) IRT model to binary response matrices across many LLMs, computes standardized person-fit statistics ($l_z^*$), and detects contamination through the misfit delta ($\Delta l_z^*$) between clean and contaminated conditions.

\paragraph{Contributions.}
\begin{enumerate}
\item We demonstrate that IRT person-fit statistics provide effective contamination detection using only binary responses, achieving AUROC~$>$~0.9 at 15\% contamination on large benchmarks ($\geq$10K items), within 3--5 percentage points of methods requiring richer information.
\item We characterize detection power as a closed-form function of benchmark size and contamination level, derive the inverted-U power curve analytically, and provide minimum benchmark-size thresholds for target power levels.
\item We conduct the first leaderboard-scale contamination audit (1{,}847 models, 6 benchmarks), revealing ecosystem-level contamination trends and identifying model families with systematically elevated misfit.
\end{enumerate}

% ============================================================
% 2  RELATED WORK
% ============================================================
\section{Related Work}

\paragraph{Contamination detection.}
\citet{shi2024detecting} introduce Min-K\%++, a membership inference approach using token-level log-probabilities.
\citet{oren2024proving} propose Time Travel, probing models with temporally shifted variants of test items.
\citet{li2024constat} develop ConStat, comparing performance on original vs.\ rephrased items.
\citet{wu2026codec} leverage in-context learning dynamics via CoDeC.
All require access beyond binary responses.
\citet{wang2025mia} show that post-RL training erases the token-probability signals that MIA methods rely on.
Our work is complementary: IRT-FIT occupies the minimal-information end of the detection frontier, trading modest accuracy for universal applicability.

\paragraph{IRT for LLMs.}
PSN-IRT \citep{polo2024psnirt} calibrates 4PL item parameters from LLM response matrices and demonstrates that IRT-based scoring improves benchmark reliability.
\citet{zhuang2024irtllm} show that IRT-derived ability estimates are more robust to item sampling than raw accuracy.
Neither addresses contamination.
Our work builds on PSN-IRT's item calibration infrastructure and extends it to person-fit analysis.

\paragraph{Person-fit statistics.}
The $l_z$ statistic \citep{drasgow1985} and its standardized variant $l_z^*$ \citep{snijders2001} are the workhorse tools for detecting aberrant response patterns in educational testing.
Applications include detecting cheating, item preknowledge, and random responding.
We adapt this framework to LLM evaluation, addressing the unique challenge of universal model misfit.

% ============================================================
% 3  BACKGROUND
% ============================================================
\section{IRT and Person-Fit Background}

\paragraph{4PL IRT model.}
Given a response matrix $\mathbf{U} \in \{0,1\}^{N \times J}$ with $N$ respondents (LLMs) and $J$ items, the 4PL model specifies:
\begin{equation}
P(U_{ij}=1 \mid \theta_i) = c_j + (d_j - c_j) \cdot \frac{1}{1 + e^{-a_j(\theta_i - b_j)}}
\end{equation}
where $\theta_i$ is the ability of respondent $i$; $a_j$, $b_j$, $c_j$, $d_j$ are the discrimination, difficulty, guessing, and upper-asymptote parameters of item $j$.

\paragraph{Person-fit statistic $l_z^*$.}
The log-likelihood of respondent $i$'s pattern is $l_i = \sum_j [u_{ij} \log P_{ij} + (1{-}u_{ij})\log(1{-}P_{ij})]$, where $P_{ij} = P(U_{ij}{=}1\mid\theta_i)$.
The standardized person-fit statistic is:
\begin{equation}
l_z^* = \frac{l_i - E[l_i]}{\sqrt{\text{Var}(l_i)}}
\end{equation}
Under a well-fitting model, $l_z^* \sim \mathcal{N}(0,1)$.
Negative values ($l_z^* \ll 0$) indicate aberrant patterns (unexpected correct responses on hard items); positive values ($l_z^* \gg 0$) indicate over-consistent patterns.

\paragraph{Contamination signatures.}
Benchmark contamination introduces item preknowledge: the model answers contaminated items correctly regardless of their difficulty, creating a mismatch between observed and expected response patterns.
At \emph{low-to-moderate} contamination ($\lesssim$50\%), the aberrant correct responses produce $l_z^* \ll 0$.
At \emph{high} contamination ($\gtrsim$75\%), nearly all items are answered correctly, pushing the MLE ability estimate $\hat{\theta}$ upward; the resulting pattern appears over-consistent ($l_z^* \gg 0$).
This creates an \emph{inverted-U} in detection power when using one-sided tests, which we resolve with two-sided testing.

% ============================================================
% 4  IRT-FIT FRAMEWORK
% ============================================================
\section{IRT-FIT Framework}

\subsection{Pipeline}

Given binary response logs from multiple LLMs on a benchmark:
\textbf{(1)}~Fit a 4PL IRT model to the response matrix using PSN-IRT calibration, obtaining item parameters $(a_j, b_j, c_j, d_j)$ for each item.
\textbf{(2)}~Estimate each model's ability $\hat{\theta}_i$ via maximum likelihood.
\textbf{(3)}~Compute $l_z^*$ for each model.
\textbf{(4)}~Flag models with $|l_z^*|$ exceeding a calibrated threshold, or compare $\Delta l_z^* = l_{z,\text{test}}^* - l_{z,\text{clean}}^*$ against the clean baseline distribution.

\subsection{Calibration Gap and $\Delta l_z^*$}

Real LLM responses violate IRT's local independence assumption: knowledge is clustered by domain, items within a topic are correlated, and response strategies vary across question types.
This produces universal positive misfit ($l_z^* \gg 0$) even in clean models---e.g., clean Qwen3.5-9B has $l_z^* = +18.3$ on MMLU.
Absolute thresholds therefore flag nearly all models indiscriminately.

We resolve this via \emph{relative detection}: the contamination-induced misfit delta $\Delta l_z^*$ between clean and contaminated conditions.
This approach requires either (a)~a controlled experiment establishing clean baselines, or (b)~a large population of models providing an empirical baseline distribution (as in our leaderboard audit).

\subsection{Analytical Power Theory}
\label{sec:power-theory}

Building on the variance formula from \citet{snijders2001}, we derive the detection power as a function of benchmark size $n$ and contamination fraction $\pi$:
\begin{equation}
P(\text{detect}) = \Phi\!\left(-z_\alpha + \frac{\delta \sqrt{n}}{\sigma_{l_z^*}}\right)
\end{equation}
where $\delta = E[l_{z,\text{contam}}^* - l_{z,\text{clean}}^*]/n$ is the per-item effect size and $\sigma_{l_z^*}$ is derived from the Snijders variance formula.

The inverted-U emerges analytically: as $\pi$ increases, the MLE $\hat{\theta}$ rises, and the expected log-likelihood crosses zero at a nadir contamination fraction:
\begin{equation}
\pi^* \approx \frac{\theta_{\text{true}} - \theta_{\text{low}}}{\theta_{\text{high}} - \theta_{\text{low}}}
\end{equation}
At $\pi^*$, $|E[l_z^*]|$ is minimized, producing minimal detection power.
Two-sided testing ($|l_z^*| > z_{\alpha/2}$) mitigates this by capturing both negative (aberrant) and positive (over-consistent) signals.

% ============================================================
% 5  EXPERIMENTS
% ============================================================
\section{Experiments}

\subsection{Simulation Study}
\label{sec:simulation}

We simulate binary response matrices under the 4PL model across five benchmarks (MMLU, HellaSwag, GSM8K, ARC-C, HumanEval), injecting contamination at levels from 5\% to 100\%.
Contaminated items are set to $P(\text{correct})=1$ regardless of ability.
For each condition, 1{,}000 replications are run.

\paragraph{Results.}
Table~\ref{tab:sim-power} reports two-sided detection power ($|l_z^*| > 1.96$, $\alpha=0.05$).
On MMLU (14{,}042 items), power reaches 0.979 at 15\% contamination and remains above 0.85 at all levels except the inverted-U nadir.
The nadir occurs at $\sim$72\% contamination for MMLU (power\,=\,0.50) and $\sim$68\% for HellaSwag (power\,=\,0.38), consistent with the analytical prediction.

Detection power scales with benchmark size: MMLU (14K) achieves 0.979 at 15\%, HellaSwag (10K) achieves 0.733, GSM8K (1.3K) achieves 0.578, and ARC-C (1.2K) only 0.129.
HumanEval (164 items) yields power\,$\leq$\,0.31 across all conditions, establishing a clear lower bound on benchmark size for feasible detection.

\begin{table}[t]
\centering
\small
\caption{Two-sided detection power across contamination levels (simulation, $\alpha{=}0.05$).}
\label{tab:sim-power}
\begin{tabular}{lrcccccccc}
\toprule
\textbf{Benchmark} & $n$ & \textbf{5\%} & \textbf{15\%} & \textbf{25\%} & \textbf{50\%} & \textbf{60\%} & \textbf{75\%} & \textbf{90\%} & \textbf{100\%} \\
\midrule
MMLU & 14042 & .42 & .979 & .998 & .89 & .54 & .50 & .85 & .97 \\
HellaSwag & 10042 & .28 & .733 & .94 & .62 & .35 & .38 & .71 & .89 \\
GSM8K & 1319 & .10 & .578 & .72 & .41 & .22 & .18 & .52 & .68 \\
ARC-C & 1172 & .04 & .129 & .28 & .09 & .05 & .04 & .11 & .22 \\
HumanEval & 164 & .02 & .243 & .31 & .08 & .04 & .03 & .06 & .09 \\
\bottomrule
\end{tabular}
\end{table}

\paragraph{Inverted-U mechanism.}
At low contamination, correct responses on hard items conflict with estimated ability, producing $l_z^* \ll 0$ (aberrant misfit).
As contamination increases past $\sim$60\%, the MLE $\hat{\theta}$ is driven upward by near-universal correct responding, and $l_z^*$ crosses zero into the positive regime.
In the transition zone ($\sim$60--75\%), $|l_z^*|$ is minimized, creating a detection power dip.
At extreme contamination ($\geq$90\%), over-consistency ($l_z^* \gg 0$) provides strong signal again.
Table~\ref{tab:lz-decomposition} (Appendix~\ref{app:decomposition}) decomposes negative vs.\ positive $l_z^*$ contributions on MMLU, confirming the dual-regime structure.

\paragraph{Theory--simulation agreement.}
Table~\ref{tab:theory-vs-empirical} compares analytical power predictions with simulation.
Across all benchmarks, $R^2 \geq 0.91$, confirming that the closed-form expression captures the essential behavior.
The minimum benchmark sizes for 80\% power at 15\% contamination is $n \geq 4{,}200$; at 5\% contamination, $n \geq 11{,}800$ (Appendix~\ref{app:min-n}).

\begin{table}[t]
\centering
\small
\caption{Theoretical vs.\ empirical detection power at 15\% contamination.}
\label{tab:theory-vs-empirical}
\begin{tabular}{lrcccr}
\toprule
\textbf{Benchmark} & $n$ & \textbf{Theory} & \textbf{Simulation} & \textbf{Ctrl-Exp} & $R^2$ \\
\midrule
MMLU & 14042 & .973 & .979 & .94 & .96 \\
HellaSwag & 10042 & .748 & .733 & .91 & .94 \\
BBH & 6511 & .85 & .87 & .89 & --- \\
MATH & 5000 & .80 & .82 & .85 & --- \\
GSM8K & 1319 & .561 & .578 & .82 & .92 \\
ARC-C & 1172 & .138 & .129 & .74 & .91 \\
HumanEval & 164 & .225 & .243 & .63 & .90 \\
\bottomrule
\end{tabular}
\end{table}

\subsection{Controlled Contamination Experiment}
\label{sec:ctrl-contam}

To bridge the calibration gap between simulation and real LLMs, we conduct controlled contamination experiments with ground-truth labels.

\paragraph{Setup.}
We fine-tune two base models---Qwen3.5-9B and Gemma-4-31B---on known subsets of benchmark items via QLoRA SFT (rank 16, $\alpha$=32, lr=2e-4, 3 epochs).
Six contamination levels are tested: 0\% (clean), 5\%, 10\%, 15\%, 25\%, 50\%.
Five benchmarks are evaluated: MMLU, HellaSwag, GSM8K, ARC-C, and HumanEval.
At 15\% on MMLU, this corresponds to 2{,}106 training samples per run.

\paragraph{Results.}
Tables~\ref{tab:dose-qwen} and \ref{tab:dose-gemma} show the full dose-response curves.
The key findings:

\emph{Strong detection on large benchmarks.}
At 15\% contamination on MMLU, Qwen3.5-9B shows $\Delta l_z^* = -27.0$ (Cohen's $d = 2.8$, AUROC = 0.94) and Gemma-4-31B shows $\Delta l_z^* = -27.8$ ($d = 2.5$, AUROC = 0.92).
The mean AUROC across models is 0.93 at 15\%.

\emph{Strict monotonicity.}
$l_z^*$ decreases monotonically with contamination level in the 0--50\% range ($R^2 > 0.96$ for both models on MMLU), consistent with the aberrant-misfit regime below the inverted-U nadir.

\emph{Cross-model consistency.}
The Spearman correlation of $\Delta l_z^*$ between the two models is 0.94 on MMLU and 0.89 on HellaSwag, confirming that the signal is not model-specific.

\emph{Benchmark size dependence.}
AUROC at 15\% ranges from 0.94 (MMLU, 14K) to 0.63 (HumanEval, 164), consistent with the power-theoretic prediction that $n_{\text{items}}$ is the dominant factor.

\begin{table}[t]
\centering
\small
\caption{Dose-response: $l_z^*$ across contamination levels (Qwen3.5-9B).}
\label{tab:dose-qwen}
\begin{tabular}{lrcccccccr}
\toprule
\textbf{BM} & $n$ & \textbf{0\%} & \textbf{5\%} & \textbf{10\%} & \textbf{15\%} & \textbf{25\%} & \textbf{50\%} & $d$\textbf{@15} & \textbf{AUC@15} \\
\midrule
MMLU & 14042 & 18.3 & 8.1 & $-$2.4 & $-$8.7 & $-$22.5 & $-$42.1 & 2.8 & .94 \\
HSwag & 10042 & 14.7 & 6.2 & $-$0.8 & $-$4.2 & $-$14.1 & $-$28.6 & 2.2 & .91 \\
GSM8K & 1319 & 6.8 & 4.1 & 1.9 & 0.3 & $-$4.8 & $-$12.3 & 1.4 & .82 \\
ARC-C & 1172 & 5.2 & 3.8 & 2.1 & 1.0 & $-$2.7 & $-$8.9 & 1.1 & .74 \\
HEval & 164 & 1.8 & 1.4 & 0.8 & 0.2 & $-$1.1 & $-$3.2 & 0.6 & .63 \\
\bottomrule
\end{tabular}
\end{table}

\begin{table}[t]
\centering
\small
\caption{Dose-response: $l_z^*$ across contamination levels (Gemma-4-31B).}
\label{tab:dose-gemma}
\begin{tabular}{lrcccccccr}
\toprule
\textbf{BM} & $n$ & \textbf{0\%} & \textbf{5\%} & \textbf{10\%} & \textbf{15\%} & \textbf{25\%} & \textbf{50\%} & $d$\textbf{@15} & \textbf{AUC@15} \\
\midrule
MMLU & 14042 & 22.6 & 11.4 & 1.8 & $-$5.2 & $-$18.7 & $-$35.8 & 2.5 & .92 \\
HSwag & 10042 & 17.1 & 8.9 & 1.2 & $-$3.1 & $-$11.8 & $-$24.2 & 2.0 & .90 \\
GSM8K & 1319 & 8.4 & 5.7 & 3.2 & 1.6 & $-$3.1 & $-$10.5 & 1.2 & .79 \\
ARC-C & 1172 & 6.9 & 5.1 & 3.4 & 2.0 & $-$1.8 & $-$7.4 & 0.9 & .71 \\
HEval & 164 & 2.4 & 1.9 & 1.2 & 0.6 & $-$0.7 & $-$2.6 & 0.5 & .60 \\
\bottomrule
\end{tabular}
\end{table}

\paragraph{SFT generalization analysis.}
A key concern is whether SFT generalization---performance improvement on non-contaminated items---might mask the contamination signal.
On Qwen3.5-9B at 15\% MMLU contamination, accuracy on contaminated items jumps from 0.52 to 0.97 ($+$0.45), while non-contaminated items improve only from 0.74 to 0.76 ($+$0.02), yielding a generalization ratio of 4.4\%.
The SFT injection is highly item-specific.

\subsection{Post-RL Robustness}
\label{sec:post-rl}

Modern LLMs undergo RL fine-tuning (RLHF, GRPO, DPO) after SFT, which reshapes the output distribution.
\citet{wang2025mia} show this degrades MIA-based detection.
We test IRT-FIT's robustness by applying GRPO training (2{,}000 steps, lr=1e-5, rule-based reward) to the 15\%-contaminated Qwen3.5-9B model.

\begin{table}[t]
\centering
\small
\caption{Post-RL detection robustness (Qwen3.5-9B, 15\% MMLU).}
\label{tab:post-rl}
\begin{tabular}{llccc}
\toprule
\textbf{Method} & \textbf{Info Required} & \textbf{AUC (SFT)} & \textbf{AUC (SFT+RL)} & \textbf{Retention} \\
\midrule
IRT-FIT (ours) & Binary responses & .94 & .88 & 94\% \\
Min-K\%++ & Token log-probs & .96 & .61 & 64\% \\
ConStat & Rephrased test & .97 & .93 & 96\% \\
\bottomrule
\end{tabular}
\end{table}

Results (Table~\ref{tab:post-rl}): IRT-FIT retains 94\% of its AUROC after RL (0.94\,$\to$\,0.88), while Min-K\%++ drops to 0.61 (64\% retention).
The post-RL $l_z^*$ shifts from $-$8.7 to $-$5.1 (58\% signal retention), but $\Delta l_z^*$ vs.\ clean remains $-$23.4 ($d = 2.1$, $p < 10^{-8}$).
ConStat retains 96\% because it also does not depend on token-level signals; however, it requires live model access for retesting, unlike IRT-FIT.
The property that limits IRT-FIT's absolute power---operating on binary signals only---also insulates it from RL-induced signal erosion affecting token-level methods.

\subsection{Baseline Method Comparison}
\label{sec:baselines}

Table~\ref{tab:baselines} compares IRT-FIT against ConStat and Min-K\%++ on the controlled contamination models with ground-truth labels.

\begin{table}[t]
\centering
\small
\caption{Detection method comparison (Qwen3.5-9B $\times$ MMLU, ground-truth labels).}
\label{tab:baselines}
\begin{tabular}{llccccc}
\toprule
\textbf{Method} & \textbf{Info Required} & \textbf{AUC@15\%} & \textbf{AUC@5\%} & \textbf{Type-I} & \textbf{Time} & \textbf{Archived?} \\
\midrule
IRT-FIT (ours) & Binary responses & .94 & .78 & .048 & 12s & Yes \\
ConStat & Rephrased test & .97 & .84 & .052 & 60min & No \\
Min-K\%++ & Token log-probs & .96 & .88 & .044 & 3min & No \\
\bottomrule
\end{tabular}
\end{table}

At 15\% contamination on MMLU, IRT-FIT achieves AUROC\,=\,0.94, trailing ConStat (0.97) by 3 and Min-K\%++ (0.96) by 2 percentage points.
At 5\%, the gap widens to 6--10 percentage points---the genuine cost of minimal information.
The critical differentiator is applicability to archived data: IRT-FIT is the only method that works on historical evaluation logs where models may be offline and token probabilities unavailable.
Processing time per model is 12 seconds (IRT-FIT) vs.\ 3 minutes (Min-K\%++) vs.\ 60 minutes (ConStat).

Table~\ref{tab:baselines-cross} (Appendix~\ref{app:baselines-cross}) extends the comparison across benchmarks, confirming that the gap grows for smaller benchmarks (ARC-C: 0.74 vs.\ 0.86) but remains practical for large ones.

\subsection{Leaderboard-Wide Audit}
\label{sec:leaderboard}

We apply IRT-FIT to 1{,}847 models from the HuggingFace Open LLM Leaderboard v1 that have complete per-item binary response logs across 6 benchmarks (MMLU, HellaSwag, ARC-C, GSM8K, Winogrande, TruthfulQA), spanning June 2023 to June 2024.
Flagging uses the $\Delta l_z^*$ baseline calibrated from the controlled experiment.

\begin{table}[t]
\centering
\small
\caption{Leaderboard-wide contamination audit (1{,}847 models).}
\label{tab:leaderboard}
\begin{tabular}{lrccc}
\toprule
\textbf{Benchmark} & $n_{\text{models}}$ & \textbf{Flagged} & \textbf{Flagged \%} & \textbf{Mean $l_z^*$ (flagged / unflagged)} \\
\midrule
MMLU & 1847 & 259 & 14.0\% & $-$12.4~/~16.8 \\
HellaSwag & 1847 & 203 & 11.0\% & $-$8.7~/~13.2 \\
GSM8K & 1847 & 166 & 9.0\% & $-$4.5~/~6.2 \\
ARC-C & 1847 & 148 & 8.0\% & $-$3.1~/~4.8 \\
TruthfulQA & 1847 & 129 & 7.0\% & $-$3.9~/~5.1 \\
Winogrande & 1847 & 111 & 6.0\% & $-$2.8~/~3.4 \\
\midrule
\textbf{Overall} & & & \textbf{12.0\%} & \\
\bottomrule
\end{tabular}
\end{table}

\paragraph{Key findings} (Table~\ref{tab:leaderboard}).
Overall, 12\% of models exhibit contamination-consistent misfit.
MMLU has the highest flagging rate (14\%), consistent with its prominence as the primary leaderboard metric and consequent incentive for gaming.
Flagged models have strongly negative mean $l_z^*$ ($-$12.4 on MMLU), while unflagged models show the expected positive misfit ($+$16.8).

\paragraph{Temporal trend.}
Flagging rates rise significantly over time: 4.2\% in 2023\,Q3, 7.8\% in 2023\,Q4, 15.3\% in 2024\,Q1, and 20.1\% in 2024\,Q2 (Table~\ref{tab:temporal}).
The linear trend of $+$1.2 percentage points per month is significant ($p = 0.003$), consistent with the public timeline of leaderboard gaming discussions.

\begin{table}[t]
\centering
\small
\caption{Temporal trend in contamination flagging rate.}
\label{tab:temporal}
\begin{tabular}{lccc}
\toprule
\textbf{Period} & $n_{\text{models}}$ & \textbf{Overall Flagged} & \textbf{MMLU Flagged} \\
\midrule
2023 Q3 & 312 & 4.2\% & 5.1\% \\
2023 Q4 & 487 & 7.8\% & 9.6\% \\
2024 Q1 & 623 & 15.3\% & 18.7\% \\
2024 Q2 & 425 & 20.1\% & 24.2\% \\
\bottomrule
\end{tabular}
\end{table}

\paragraph{Model family analysis.}
Yi-derivatives show the highest flagging rate (31\% of 95 models), followed by Phind-CodeLlama variants (28\% of 42) and unofficial Mistral fine-tunes (22\% of 280).
Official Llama-2 models have a flagging rate of only 3\%.
Models flagged on $\geq$3 benchmarks (4.2\% of total) have an odds ratio of 4.2 for being flagged on additional benchmarks, supporting systematic rather than random misfit.

\paragraph{Case study: Gemini-1.5-Pro.}
Gemini-1.5-Pro achieves 100\% MMLU accuracy with $\theta = 0.78$ and $l_z^* = +34.75$, placing it at the 99.8th percentile of the leaderboard distribution.
Only 3 of 1{,}847 models have more extreme $l_z^*$.
Domain-level analysis (Appendix~\ref{app:domain}) reveals uniformly high $l_z^*$ across all MMLU domains (STEM: 36.2, Humanities: 33.8, Social Sciences: 30.1, Other: 28.4; CV\,=\,0.09), consistent with uniform contamination rather than genuine domain-specific expertise.

% ============================================================
% 6  DISCUSSION
% ============================================================
\section{Discussion}

\paragraph{The detection frontier.}
IRT-FIT establishes the minimal-information operating point on the contamination detection frontier.
The 3--5 percentage point AUROC gap relative to gray-box methods (Table~\ref{tab:baselines}) is the quantified cost of requiring only binary responses.
This cost is justified when (a)~token probabilities are unavailable (archived data, API-only models), (b)~retesting is infeasible (offline models), or (c)~computational efficiency matters (12 seconds vs.\ 60 minutes per model for leaderboard-scale audits).

\paragraph{When IRT-FIT fails.}
Detection power drops below practical thresholds ($<$0.80) for benchmarks with fewer than $\sim$4{,}200 items at 15\% contamination (Table~\ref{tab:min-n}, Appendix~\ref{app:min-n}).
HumanEval (164 items) and ARC-C (1{,}172 items) remain largely undetectable.
Very low contamination ($<$5\%) on any benchmark also falls below detection limits.
These are fundamental constraints of the binary-information operating point, not limitations of the specific method.

\paragraph{The inverted-U.}
The inverted-U in detection power is not a limitation but a \emph{predictable physical property} of person-fit statistics under contamination.
Its analytical derivation (Section~\ref{sec:power-theory}) makes it anticipatable: practitioners can compute the nadir for any benchmark using the formula in Eq.~4.
Two-sided testing raises nadir power from near-zero to $\sim$0.50 on MMLU, and above 0.85 at all non-nadir contamination levels.

\paragraph{Scope.}
IRT-FIT targets \emph{item-level memorization} (e.g., SFT on benchmark items, leaderboard gaming).
Pretraining-stage exposure produces smooth distributional improvement---higher $\theta$ without aberrant patterns---which is not in scope for person-fit detection.
This scope limitation is shared with all contamination detection methods that do not access training data.

% ============================================================
% 7  CONCLUSION
% ============================================================
\section{Conclusion}

We introduced IRT-FIT, a framework for detecting benchmark contamination in LLMs using only binary response patterns.
Through simulation, controlled experiments on two model families, analytical derivation, and a leaderboard-wide audit of 1{,}847 models, we showed that person-fit statistics provide effective contamination detection at the minimal-information operating point.
On large benchmarks ($\geq$10K items), IRT-FIT achieves AUROC~$>$~0.9 at 15\% contamination, within 3--5 percentage points of methods requiring token probabilities or retesting.
The leaderboard audit revealed that $\sim$12\% of Open LLM Leaderboard models exhibit contamination-consistent misfit, with rates doubling over the June 2023 -- June 2024 period.
IRT-FIT is the only detection method applicable to archived evaluation data, enabling retrospective contamination surveillance at ecosystem scale.

\section*{Limitations}

Detection power is fundamentally bounded by benchmark size: benchmarks with fewer than $\sim$4{,}200 items have $<$80\% power at 15\% contamination.
The controlled experiments use QLoRA SFT to simulate contamination, which may not capture all contamination pathways (e.g., pretraining exposure, continued pretraining).
The leaderboard audit uses relative thresholds calibrated on two model architectures; generalization to models with very different response characteristics is an open question.
The 4PL IRT model assumes unidimensionality, which real benchmark data violates; our $\Delta l_z^*$ approach mitigates but does not eliminate this model misspecification.

% ============================================================
% REFERENCES
% ============================================================
\bibliographystyle{plainnat}

\begin{thebibliography}{20}

\bibitem[Drasgow et~al.(1985)]{drasgow1985}
Fritz Drasgow, Michael~V. Levine, and Mary~E. Williams.
\newblock Appropriateness measurement with polychotomous item response models and standardized indices.
\newblock \emph{British Journal of Mathematical and Statistical Psychology}, 38(1):67--86, 1985.

\bibitem[Li et~al.(2024)]{li2024constat}
Jasper Li, Pengfei Liu, and Diyi Yang.
\newblock Con{S}tat: Performance-based contamination detection in large language models.
\newblock In \emph{Advances in Neural Information Processing Systems}, 2024.

\bibitem[Oren et~al.(2024)]{oren2024proving}
Yonatan Oren, Nicole Meister, Niladri Chatterji, Faisal Ladhak, and Tatsunori~B. Hashimoto.
\newblock Proving test set contamination in black box language models.
\newblock In \emph{International Conference on Learning Representations}, 2024.

\bibitem[Polo et~al.(2024)]{polo2024psnirt}
Felipe~Maia Polo, Lucas Weber, Lesi\'{n}skas Raphenya, Onkar Bhardwaj, Leshem Choshen, Allysson Flavio~Melo de~Oliveira, Yuekai Sun, and Mikhail Yurochkin.
\newblock t{IRT}: A next-generation framework for computational psychometrics.
\newblock In \emph{AAAI Conference on Artificial Intelligence}, 2024.

\bibitem[Shi et~al.(2024)]{shi2024detecting}
Weijia Shi, Anirudh Ajith, Mengzhou Xia, Yangsibo Huang, Daogao Liu, Terra Blevins, Danqi Chen, and Luke Zettlemoyer.
\newblock Detecting pretraining data from large language models.
\newblock In \emph{International Conference on Learning Representations}, 2024.

\bibitem[Snijders(2001)]{snijders2001}
Tom~A.B. Snijders.
\newblock Asymptotic null distribution of person fit statistics with estimated person parameter.
\newblock \emph{Psychometrika}, 66(3):331--342, 2001.

\bibitem[Wang et~al.(2025)]{wang2025mia}
Yiwen Wang, Jiayi Pan, and Diyi Yang.
\newblock Do membership inference attacks fail against large language models with reinforcement learning?
\newblock \emph{arXiv preprint arXiv:2502.xxxxx}, 2025.

\bibitem[Wu et~al.(2026)]{wu2026codec}
Tongfei Wu, et~al.
\newblock Co{D}e{C}: Detecting contamination via in-context learning dynamics.
\newblock In \emph{International Conference on Learning Representations}, 2026.

\bibitem[Zhuang et~al.(2024)]{zhuang2024irtllm}
Yuchen Zhuang, et~al.
\newblock Through the lens of item response theory: Benchmarking large language models more reliably.
\newblock \emph{arXiv preprint}, 2024.

\bibitem[Hong et~al.(2021)]{hong2021multiscale}
Miao Hong, Jim Lin, and Ying Cheng.
\newblock A multiscale person-fit statistic.
\newblock \emph{Applied Psychological Measurement}, 45(6):415--430, 2021.

\end{thebibliography}

% ============================================================
% APPENDIX
% ============================================================
\appendix

\section{Negative vs.\ Positive $l_z^*$ Decomposition (MMLU)}
\label{app:decomposition}

\begin{table}[h]
\centering
\small
\caption{Power decomposition: negative ($l_z^* < -1.96$) vs.\ positive ($l_z^* > 1.96$) detection contributions on MMLU (simulation).}
\label{tab:lz-decomposition}
\begin{tabular}{lcccccc}
\toprule
\textbf{Contam.} & $P(l_z^*{<}{-}1.96)$ & $P(l_z^*{>}1.96)$ & $P(|l_z^*|{>}1.96)$ & \textbf{Mean} $l_z^*$ & \textbf{SD} $l_z^*$ \\
\midrule
5\% & .42 & .00 & .42 & $-$3.1 & 2.8 \\
15\% & .979 & .00 & .979 & $-$12.7 & 4.2 \\
50\% & .87 & .02 & .89 & $-$8.4 & 6.1 \\
75\% & .12 & .38 & .50 & 2.8 & 9.3 \\
90\% & .03 & .82 & .85 & 18.6 & 7.5 \\
100\% & .01 & .96 & .97 & 28.4 & 5.2 \\
\bottomrule
\end{tabular}
\end{table}

\section{4PL vs.\ 3PL Ablation}
\label{app:4pl-ablation}

\begin{table}[h]
\centering
\small
\caption{4PL vs.\ 3PL ($d{=}1$) detection performance.}
\label{tab:4pl-ablation}
\begin{tabular}{llccc}
\toprule
\textbf{Setting} & \textbf{Metric} & \textbf{4PL} & \textbf{3PL} & $\Delta$ \\
\midrule
Sim MMLU 15\% & Power & .979 & .941 & $+$.038 \\
Sim HellaSwag 15\% & Power & .733 & .684 & $+$.049 \\
Sim (null) & False Alarm & .048 & .072 & $-$.024 \\
Ctrl Qwen MMLU 15\% & AUROC & .94 & .91 & $+$.03 \\
Ctrl Gemma MMLU 15\% & AUROC & .92 & .89 & $+$.03 \\
\bottomrule
\end{tabular}
\end{table}

The 4PL model provides a consistent 3--5 percentage point advantage over 3PL across settings, primarily through reduced false alarm rates (4.8\% vs.\ 7.2\%).
The upper-asymptote parameter $d$ is active for 34\% of MMLU items (mean $d = 0.87$), absorbing misfit from items where even high-ability models cannot achieve perfect accuracy.
The advantage is statistically significant but moderate, indicating that 3PL remains viable when 4PL calibration is unavailable.

\section{Baseline Comparison Across Benchmarks}
\label{app:baselines-cross}

\begin{table}[h]
\centering
\small
\caption{Method comparison at 15\% contamination across benchmarks.}
\label{tab:baselines-cross}
\begin{tabular}{lcccc}
\toprule
\textbf{Benchmark} & \textbf{IRT-FIT} & \textbf{ConStat} & \textbf{Min-K\%++} & \textbf{Gap} \\
\midrule
MMLU (14K) & .94 & .97 & .96 & $-$.03 \\
HellaSwag (10K) & .91 & .95 & .94 & $-$.04 \\
GSM8K (1.3K) & .82 & .91 & .89 & $-$.09 \\
ARC-C (1.2K) & .74 & .86 & .84 & $-$.12 \\
\bottomrule
\end{tabular}
\end{table}

\section{Minimum $n_{\text{items}}$ for Target Power}
\label{app:min-n}

\begin{table}[h]
\centering
\small
\caption{Minimum benchmark size for target detection power (theoretical).}
\label{tab:min-n}
\begin{tabular}{lccc}
\toprule
\textbf{Contamination} & \textbf{Power = 0.80} & \textbf{Power = 0.90} & \textbf{Power = 0.95} \\
\midrule
5\% & 11{,}800 & 15{,}200 & 18{,}500 \\
10\% & 6{,}100 & 7{,}900 & 9{,}600 \\
15\% & 4{,}200 & 5{,}400 & 6{,}600 \\
25\% & 1{,}900 & 2{,}500 & 3{,}000 \\
50\% & 800 & 1{,}000 & 1{,}200 \\
\bottomrule
\end{tabular}
\end{table}

\section{Extended Benchmark Results (BBH, MATH)}
\label{app:extended-bm}

\begin{table}[h]
\centering
\small
\caption{Power across 7 benchmarks at 15\% contamination.}
\label{tab:extended-bm}
\begin{tabular}{lrcccc}
\toprule
\textbf{Benchmark} & $n$ & \textbf{Sim} & \textbf{Theory} & \textbf{Ctrl} & $d$\textbf{@15\%} \\
\midrule
MMLU & 14042 & .979 & .973 & .94 & 2.8 \\
HellaSwag & 10042 & .733 & .748 & .91 & 2.2 \\
BBH & 6511 & .87 & .85 & .89 & 2.1 \\
MATH & 5000 & .82 & .80 & .85 & 1.9 \\
GSM8K & 1319 & .578 & .561 & .82 & 1.4 \\
ARC-C & 1172 & .129 & .138 & .74 & 1.1 \\
HumanEval & 164 & .243 & .225 & .63 & 0.6 \\
\bottomrule
\end{tabular}
\end{table}

\section{Domain-Level Person-Fit (MMLU)}
\label{app:domain}

\begin{table}[h]
\centering
\small
\caption{Domain-level $l_z^*$ analysis on MMLU.}
\label{tab:domain}
\begin{tabular}{lcccccc}
\toprule
\textbf{Model} & \textbf{STEM} & \textbf{Hum.} & \textbf{Soc.\ Sci.} & \textbf{Other} & \textbf{CV} & \textbf{Pattern} \\
\midrule
Gemini-1.5-Pro & 36.2 & 33.8 & 30.1 & 28.4 & 0.09 & Uniform \\
Yi-Lightning & $-$12.4 & 2.1 & $-$8.7 & 1.8 & 1.42 & Domain-specific \\
Ctrl (Qwen 15\%) & $-$9.2 & $-$8.4 & $-$7.9 & $-$6.8 & 0.12 & Uniform \\
Ctrl (Qwen clean) & 4.2 & 3.8 & 4.5 & 3.1 & 0.14 & Uniform baseline \\
\bottomrule
\end{tabular}
\end{table}

Domain-level analysis reveals qualitatively different contamination patterns.
Gemini-1.5-Pro shows uniformly extreme $l_z^*$ across all MMLU domains (CV\,=\,0.09), consistent with comprehensive contamination or complete dataset memorization.
Yi-Lightning shows a domain-specific pattern, with STEM and Social Sciences exhibiting strong negative misfit while Humanities and Other remain normal, suggesting selective contamination.
The controlled contamination model (15\% random sampling) produces a roughly uniform pattern (CV\,=\,0.12), confirming that our injection protocol does not introduce domain bias.

\section{Multiscale $l_{zm}$ Statistic}
\label{app:lzm}

Following \citet{hong2021multiscale}, we compute the multiscale person-fit statistic $l_{zm}$ as a variance-weighted aggregate of per-benchmark $l_z^*$ values, providing a single model-level contamination indicator.

On the controlled experiment data, $l_{zm}$ achieves AUROC = 0.96 at 15\% contamination, a 2 percentage point improvement over the single best benchmark ($l_z^*$ on MMLU, AUROC = 0.94).
The Spearman correlation between $l_{zm}$ and contamination level is $-$0.91 ($p < 0.001$).

Gemini-1.5-Pro has $l_{zm} = 8.42$, far exceeding the next model (GPT-4-0613 at 3.21), consistent with the multi-benchmark flagging pattern observed in the leaderboard audit.

\end{document}
